\section{Negative equality: a global weighted logarithm inequality}
\label{sec:weighted}
The local result leaves only fixed points with $\mu>1/20$.
We now exclude these except in a short interval near the coloring
channel.  The scalar inequality in this section is certified by the
integer interval calculation in Appendix~\ref{int:appendix}, with its
complete source in Section~\ref{cert:interval}.

\begin{proposition}[Global exclusion away from coloring]
\label{wei:theorem}\label{wei:exclusion}
At $\lambda=-a$, $d=a^{-2}$, the uniform posterior law is the only
permutation-invariant fixed point in either of the following ranges:
\begin{enumerate}
 \item regular offspring, with integer $d$ and $0<a\le.448$;
 \item Poisson offspring, with $0<a\le.477$.
\end{enumerate}
In particular the regular assertion covers every integer $d\ge5$.
\end{proposition}

For a positive likelihood vector $m=1+v$ define
\begin{equation}\label{wei:functions}
 G=(m_0m_1m_2)^{1/3},\quad T=1-G,\quad
 u_i=\log m_i-\av{\log m},\quad
 W=G\av{u^2},\quad Y=G\av{u^3}.
\end{equation}
Here $\av{\cdot}$ is the average of the three coordinates.  Edge
subscripts mean evaluation on $1-av$, and we use
\[
 \widehat T=dT_e,\qquad\widehat W=dW_e,\qquad\widehat Y=dY_e.
\]
The functions $T,W,Y$ extend continuously to the closed simplex, with
$T=1$ and $W=Y=0$ whenever any coordinate vanishes.  Indeed
$G$ times any fixed power of the logarithms tends to zero as a coordinate
vanishes: for example the product $\prod_i m_i^{1/3}$ times
$\sum_i|\log m_i|^3$ is bounded by a sum of terms having a factor
$m_i^{1/3}|\log m_i|^3\to0$ and the other coordinates are bounded by
three.  This observation also proves boundedness of these extensions.

\begin{lemma}[Weighted fixed-point identities]\label{wei:identities}
Let $0<a<1/2$, $d=a^{-2}$, and let the fixed point be nontrivial.
Set
\begin{equation}\label{wei:scalars}
 \Phi(t)=\E t^{D},\qquad A=\E G_e,\qquad
 z=d(1-A),\qquad h=\frac{1-\Phi(A)}{d(1-A)},\qquad
 1-\xi=\frac{\Phi'(A)}d.
\end{equation}
Then $0<A<1$, $0<\xi<1$, and
\begin{equation}\label{wei:exact}
 \E T=h\E\widehat T,\qquad
 \E W=(1-\xi)\E\widehat W,\qquad
 \E Y=(1-\xi)\E\widehat Y,\qquad
 \E\widehat T=z.
\end{equation}
\end{lemma}

\begin{proof}
Edge likelihood coordinates lie in $[1-2a,1+a]$, so $G_e>0$.
Their coordinate average is one, hence $G_e\le1$, with strict inequality
on nonuniform messages.  Nontriviality and $a>0$ thus give $0<A<1$.
The assertions about $\xi$ follow from
$\Phi'(A)/d=A^{d-1}$ in the regular case and
$\Phi'(A)/d=e^{d(A-1)}$ in the Poisson case.

For $n$ branches let $Z=\av{\prod_{j=1}^n m_{e,j}}$ be the
product likelihood normalizer, and let $g_j=G_{e,j}$.
Under the independent marginal branch law the actual output law is
tilted by $Z$, whereas its geometric mean equals $\prod_jg_j/Z$.
Their product is $\prod_jg_j$.  Centered log coordinates add,
$u_{\mathrm{out},i}=\sum_j u_{e,j,i}$, because normalization changes
all three logs by the same constant.
For each fixed coordinate $i$, permutation invariance and centering
give
\[
 \E(G_e u_{e,i})=0:
\]
the three expectations are equal and their sum is zero.
Consequently every cross term in the square or cube of
$\sum_j u_{e,j,i}$ has a vanishing first-moment factor.
Independence therefore proves, for $n\ge1$,
\[
 \E(G_{\mathrm{out}}\mid n)=A^n,\quad
 \E(W_{\mathrm{out}}\mid n)=nA^{n-1}\E W_e,\quad
 \E(Y_{\mathrm{out}}\mid n)=nA^{n-1}\E Y_e.
\]
At $n=0$ the output is uniform, so the contributions to $T,W,Y$ are
zero.  Mixing over $D$ gives~\eqref{wei:exact}.
All terms are integrable: the edge logs are uniformly bounded and the
output centered logs grow at most linearly in $n$, while both offspring
laws have finite moments of every fixed order.  No conditioning on
survival enters these identities.
\end{proof}

\begin{lemma}[Scalar bounds]\label{wei:scalarbounds}
In the regular range of Proposition~\ref{wei:exclusion}, or in its
Poisson range, one can use constants
\[
 (H,U)=\left(\frac1{16},\frac{15}{32}\right)
\]
to obtain
\begin{equation}\label{wei:hz}
 h\le\overline h:=1-\frac\xi2-H\xi^2-\frac H2\xi^3,
 \qquad z\ge\underline z:=\xi+U\xi^2.
\end{equation}
For Poisson offspring the stronger values
$(H,U)=(1/12,1/2)$ also hold.  Every nontrivial fixed point in the
claimed ranges satisfies $.03<\xi<1$.
\end{lemma}

\begin{proof}
Regular offspring in the stated range have integer $d\ge5$.
Since $1-\xi=A^{d-1}$, the finite geometric sum gives
\[
 h=\frac1d\sum_{i=0}^{d-1}(1-\xi)^{i/(d-1)}.
\]
For Poisson offspring $1-\xi=e^{-z}$ and
\[
 h=\frac{1-e^{-z}}z=\int_0^1(1-\xi)^t\,dt.
\]
In either expression the averaging law of $t$ is symmetric under
$t\mapsto1-t$, so $\E_t t=1/2$.  For $0\le t\le1$ all
nonconstant coefficients in the binomial expansion of $(1-\xi)^t$
are nonpositive.  The first three negative coefficients, after
averaging, are
\[
 \frac12,\qquad \frac{\E_t[t(1-t)]}{2},\qquad
 \frac{\E_t[t(1-t)(2-t)]}{6}
       =\frac{\E_t[t(1-t)]}{4}.
\]
The last equality uses symmetry, since the $t(1-t)$-weighted mean
of $t$ is $1/2$.  In the regular case
\[
 \E_t[t(1-t)]=\frac{d-2}{6(d-1)}\ge\frac18;
\]
in the Poisson case it is $1/6$.  Truncating after degree three gives
the bound for $h$.  It holds at $\xi=1$ as well by continuity.

For regular offspring,
\[
 z=d\{1-(1-\xi)^{1/(d-1)}\}.
\]
All coefficients of this series are nonnegative.  Its linear
coefficient is $d/(d-1)\ge1$, and its quadratic coefficient obeys
\[
 \frac{d(d-2)}{2(d-1)^2}-\frac{15}{32}
       =\frac{(d-5)(d+3)}{32(d-1)^2}\ge0.
\]
This proves the claimed lower bound.  For Poisson offspring
$z=-\log(1-\xi)\ge\xi+\xi^2/2$; the weaker values of $H,U$
are consequently valid there too.

Finally
$G_e^3=1-3a^2x-2a^3y$, so, since $d=a^{-2}$,
\begin{equation}\label{wei:That}
 \widehat T=\frac{3x+2ay}{1+G_e+G_e^2}
       \ge(1-a/3)x.
\end{equation}
For this inequality use $G_e\le1$ and $y\ge-x/2$ from
\eqref{loc:ratio}; the numerator is nonnegative.  By
Proposition~\ref{loc:gap} and $a\le.477$,
\[
 z=\E\widehat T>(1-.477/3)/20=.04205>.042.
\]
For regular $d\ge5$,
\[
 \log(1-\xi)=(d-1)\log(1-z/d)\le-(1-1/d)z\le-.8z;
\]
for Poisson, $1-\xi=e^{-z}\le e^{-.8z}$.
Since $1-e^{-t}\ge t/(1+t)$ for $t\ge0$, in either case
\[
 \xi>\frac{.0336}{1+.0336}=\frac{21}{646}>.0325>.03.
\]
This verifies the complete $\xi$ domain required below.
\end{proof}

\subsection{A pointwise inequality from exact identities}

The identities above give several functions with expectation zero.  We
combine them so that the scalar bounds on $h$ and $z$ produce a function
that is strictly positive at every message.  For a nontrivial fixed
point, choose signed rational constants $c_0,c_1,s_1$ and a positive
rational $\delta$, and put
\[
 s_0=c_0+1/3,\qquad c(\xi)=c_0+c_1\xi,\qquad
 s(\xi)=s_0+s_1\xi.
\]
The relation defining $s_0$ removes the cubic term of
$T+c_0W+s_0Y$ near the uniform message.  Indeed, expansion at $v=0$
gives
\[
 T=x-\tfrac23y+O(\|v\|^4),\qquad
 W=2x-2y+O(\|v\|^4),\qquad Y=2y+O(\|v\|^4).
\]
Its quadratic term also cancels on subtracting $d$ times its edge
evaluation, since $da^2=1$.
By Lemma~\ref{wei:identities}, the random expression
\begin{align}
 E_0={}&T-h\widehat T
       +c(\xi)[W-(1-\xi)\widehat W]
       +s(\xi)[Y-(1-\xi)\widehat Y]\nonumber\\
     &+\delta\xi(z-\widehat T)
 \label{wei:zero-mean}
\end{align}
has expectation zero.  In particular, the final term changes the
pointwise expression without changing its expectation, since
$\E\widehat T=z$.

Fix an applicable pair $(H,U)$ from Lemma~\ref{wei:scalarbounds}.
Replace $h$ in~\eqref{wei:zero-mean} by
$\overline h=1-\xi/2-H\xi^2-(H/2)\xi^3$ and $z$ by
$\underline z=\xi+U\xi^2$, and call the resulting expression $P$.
The only changes are in terms whose signs are controlled:
\begin{equation}\label{wei:comparison}
 P-E_0=(h-\overline h)\widehat T
          +\delta\xi(\underline z-z)\le0.
\end{equation}
Here $\widehat T\ge0$, $\delta>0$, and $\xi\ge0$.
The possibly negative coefficients $c(\xi)$ and $s(\xi)$ multiply
identities that remain exact.  Thus $\E P\le0$ for any nontrivial
fixed point; it remains to choose the constants so that $P>0$
pointwise.

Expanding the relaxed expression gives
\begin{equation}\label{wei:polynomial}
 P=D_0+D_1\xi+D_2\xi^2+D_3\xi^3,
\end{equation}
where
\begin{align}
 D_0={}&T-\widehat T+c_0(W-\widehat W)+s_0(Y-\widehat Y),\nonumber\\
 D_1={}&c_1(W-\widehat W)+s_1(Y-\widehat Y)
       +c_0\widehat W+s_0\widehat Y+(1/2-\delta)\widehat T,\nonumber\\
 D_2={}&c_1\widehat W+s_1\widehat Y+H\widehat T+\delta,\nonumber\\
 D_3={}&U\delta+(H/2)\widehat T.
 \label{wei:coefficients}
\end{align}
We use the following rational choices on three closed $a$ bands:
\begin{equation}\label{wei:parameters}
\begin{array}{c|rrrr|cc}
 a\text{ band}&1000c_0&1000c_1&1000s_1&1000\delta&H&U\\ \hline
 \bigl[0,.190\bigr]&-268&-60&-55&130&1/16&15/32\\
 \bigl[.190,.448\bigr]&-269&67&47&55&1/16&15/32\\
 \bigl[.448,.477\bigr]&-271&58&49&50&1/12&1/2
\end{array}
\end{equation}
The final row uses the stronger scalar bounds for Poisson offspring.
Overlapping endpoints can use either applicable row.
For purposes of the pointwise inequality, extend the hatted functions
to $a=0$ by
\[
 \widehat T=x,\qquad\widehat W=2x,\qquad\widehat Y=0.
\]
These are their continuous limits: use
$G_e^3=1-3a^2x-2a^3y$ for $\widehat T$ and
$\log(1-av_i)=-av_i+O(a^2)$ for the centered edge logarithms.

\begin{lemma}[Certified global inequality]\label{wei:positive}
For each row of~\eqref{wei:parameters}, every posterior probability
vector, every $a$ in that row's closed band, and every
$.03\le\xi\le1$, the polynomial~\eqref{wei:polynomial} satisfies
$P>0$.
\end{lemma}

\begin{proof}
Appendix~\ref{int:appendix} proves the exact interval certificate for
this inequality, with the complete source in
Section~\ref{cert:interval}.  It parametrizes sorted probability vectors
as $(1,q^{18},r^{18})/(1+q^{18}+r^{18})$, where $0\le r\le q\le1$.
In these coordinates the weighted logarithmic functions have continuous
first derivatives at zero coordinates.  The appendix derives formulas
for the hatted functions without division by $a$, so their values and
first derivatives are enclosed also at $a=0$.  It then identifies the
four expressions computed by the program with
$D_0,D_1,D_2,D_3$ above.

The certificate covers the full $(q,r,a)$ domain by closed boxes and
partitions $\xi\in[.03,1]$ at
\[
 .03,.05,.075,.10,.15,.20,.30,.40,.50,.65,.80,.90,1.
\]
On every resulting interval it certifies strictly positive lower bounds
for the four cubic Bernstein coefficients.  All enclosures, including
the derivative bounds used for midpoint acceptance and subdivision,
follow from integer
arithmetic and explicit analytic remainder bounds.  The appendix proves
coverage of every simplex boundary and band endpoint and records normal
completion of all three prescribed runs with assertions enabled.
Positive Bernstein coefficients imply $P>0$ throughout each closed
$\xi$ interval.  At the uniform message this also follows directly from
$P=\delta\xi^2+U\delta\xi^3>0$.
\end{proof}

\begin{proof}[Proof of Proposition~\ref{wei:exclusion}]
Suppose a nontrivial fixed point exists.  Lemmas
\ref{wei:identities} and~\ref{wei:scalarbounds} place its scalar $\xi$
in the certified interval.  Choose the applicable row
of~\eqref{wei:parameters}.  Equation~\eqref{wei:zero-mean} has
expectation zero, so~\eqref{wei:comparison} gives $\E P\le0$.
On the other hand, Lemma~\ref{wei:positive} gives $P>0$ pointwise.
The functions in $P$ are bounded on the closed simplex for the fixed
$a<1/2$, so $\E P>0$, a contradiction.

For the final assertion, $a=1/\sqrt d\le1/\sqrt5<.448$ whenever
$d\ge5$: the strict comparison follows from
$5(.448)^2=1.00352>1$.
\end{proof}
